\paragraph{Sensitivity design.}
For each modality $m\in\{\mathrm{text},\mathrm{audio},\mathrm{video}\}$ we apply
mean shifts of both signs ($\pm$), and for each pair $(a,b)$ we plant
\texttt{interaction\_shift} with $\lambda\in\{-1.8,+1.8\}$ (quick grid; full
grid adds intermediate magnitudes). On every cell we record
(i) power on each \emph{single} modality,
(ii) power on the concatenated \emph{pair},
(iii) the PO interaction $\hat I(a,b)$ and $p(|I|)$.

\paragraph{Signed power vs fusion lift.}
Signed power is
$\mathbb{P}(p(|I|)<\alpha \wedge \hat I>0)$ under $\lambda>0$
(and the negative analogue under $\lambda<0$).
Fusion lift is
$\mathrm{power}_{\mathrm{pair}}-\max(\mathrm{power}_a,\mathrm{power}_b)$.
The operational claim: when the plant is synergistic ($\lambda>0$),
signed power rises \emph{together with} fusion lift --- the pair beats the
better single precisely when PO-risk reports positive interaction.
When $\lambda<0$, $\hat I$ flips and the lift collapses.
Mean shifts move single-modality power with the shifted block but leave
signed interaction near the null (main effect $\neq$ interaction).

\paragraph{Why this delays delivery.}
The concept-drift T$\times$A/T$\times$V table alone cannot show this
mechanism. The sensitivity board is the unique evidence that
positive $I$, signed power, and pair-vs-single detectability are the same story.